\section{Zero side: stationary residues force negative inertia}
\label{sec:zero-side}

Straighten the critical line by
\[
F(z):=\mathscr Z\!\left(\frac12+iz\right).
\]
The symmetric rectangle \eqref{eq:symmetric-contour-rectangle} becomes
\begin{equation}
\Omega_T:=\{z:T_1\le\Re z\le T_2,\ |\Im z|\le c-\tfrac12\}.
\label{eq:zero-rectangle}
\end{equation}
Thus $\Omega_T$ is invariant under $z\mapsto\bar z$ and contains both members
of every reflected off-critical pair in the height range. Moreover
$F(\bar z)=\overline{F(z)}$ and $F$ is real on $\mathbb R$. For a test
function $h$ put $h^\#(z)=\overline{h(\bar z)}$. In zero coordinates the same
normalized contour matrix has kernel
\begin{equation}
\mathcal A_F(h_1,h_2)
=\frac1{2\pi i}\int_{\partial\Omega_T}
L^{-2}\left(-F(z)\frac{F''(z)^2}{F'(z)}\right)
 h_2(z)h_1^\#(z)\,dz.
\label{eq:zero-kernel}
\end{equation}
We write $A_F(T):=A_T$ for this original matrix. If $H$ is a sufficiently
small real-symmetric perturbation of $F$, $A_H(T)$ denotes the same contour
construction with $F$ replaced by $H$. The residue lemmas below are stated for
such a generic $H$ and are then applied to the perturbation constructed in
Section~\ref{sec:multiplicity-endgame}.

\begin{lemma}[Stationary residue blocks]
\label{lem:stationary-blocks}
If $H'(c)=0$ and $H''(c)\neq0$, then
\[
\Res_{z=c}
L^{-2}\left(-H(z)\frac{H''(z)^2}{H'(z)}\right)
=-L^{-2}H(c)H''(c)=:w_c.
\]
For $c\in\mathbb R$, the contribution is a negative scalar exactly when
$H(c)H''(c)>0$. For a nonreal conjugate pair $\rho,\bar\rho$, the residue block
is
\[
\begin{pmatrix}
0&\overline{w_\rho}\\
w_\rho&0
\end{pmatrix},
\]
which has exactly one negative eigenvalue whenever $w_\rho\neq0$.
\end{lemma}

\begin{proof}
Since
\[
H'(z)=H''(c)(z-c)+O((z-c)^2),
\]
the residue is $-L^{-2}H(c)H''(c)$. The positive scalar $L^{-2}$ does not
change its sign. The conjugate-pair block has characteristic
polynomial $\lambda^2-|w_\rho|^2$.
\end{proof}

Let $B(H)$ denote the number of real stationary points with
$H(c)H''(c)>0$, and let $P(H')$ denote the number of nonreal conjugate pairs of
stationary zeros.

\subsection{Exact packet evaluation rank}
Let $C_{\rm pkt}$ be a fixed constant with
\[
0<C_{\rm pkt}<2,
\]
chosen first in the global parameter hierarchy of Section~\ref{sec:retained},
before the remaining buffer and symbol exponents.  The upper bound is not needed
for the zero-side interpolation itself; it is reserved for the exterior
packet-support rank estimate of Proposition~\ref{prop:exterior-packet-support}.
Put
\begin{equation}
\boxed{L_{\rm pkt}:=L+C_{\rm pkt}\log L.}
\label{eq:packet-bandwidth-reserve}
\end{equation}
The central width-$L$ band is the physical zeta band; the additional
$C_{\rm pkt}\log L$ is a two-sided reserve used for packet-dimension slack and
the smooth transition caps.
Choose a real $C^\infty$ function $\phi_L$ which equals
$L_{\rm pkt}^{-1/2}$ on $[-L_{\rm pkt}/2,L_{\rm pkt}/2]$, satisfies
$|\phi_L|\le L_{\rm pkt}^{-1/2}$, and is supported in
$[-L_{\rm pkt}/2-\Delta,L_{\rm pkt}/2+\Delta]$ for one fixed $\Delta>0$. Put
\[
\psi_L(z):=\int_{\mathbb R}\phi_L(u)e^{iuz}\,du,
\qquad
\psi_{L,\tau}(z):=\psi_L(z-\tau).
\]
The reserve supplies more than the $O(T)$ dimensional slack required by the
Riemann--von Mangoldt secondary term and leaves fixed-width room for the smooth
transition caps. By Proposition~\ref{prop:trivial-outer-twist}, no arithmetic
outer translation has to be accommodated. Since
$L_{\rm pkt}=L+O(\log L)$, all normalized analytic scales below change only by
$O(\log L/L)$, and on a fixed horizontal strip the larger exponential type
costs only a fixed power of $L$.
Repeated integration by parts gives, on every fixed horizontal strip and for
every fixed $M$,
\begin{equation}
|\psi_L(z)|\ll_M L_{\rm pkt}^{1/2}(1+|\Re z|)^{-M}
 e^{(L_{\rm pkt}/2+\Delta)|\Im z|}.
\label{eq:packet-decay}
\end{equation}

\begin{lemma}[Local packet centres give surjective evaluation]
\label{lem:evaluation-rank}
Let $z_1,\dots,z_m$ be distinct complex stationary nodes and let
$I_1,\dots,I_m$ be arbitrary nonempty open intervals of admissible packet
centres. There exist $\tau_j\in I_j$ such that
\[
\det\bigl(\psi_L(z_i-\tau_j)\bigr)_{1\le i,j\le m}\ne0.
\]
Hence the packet evaluation map onto the stationary nodes is surjective.
\end{lemma}

\begin{proof}
First fix distinct $z_i$. If
\[
\sum_i c_i\psi_L(z_i-\tau)=0
\]
on a nonempty real interval in $\tau$, analyticity makes the identity global.
Fourier uniqueness then gives
\[
\phi_L(u)\sum_i c_ie^{iuz_i}=0
\]
for almost every $u$. Since $\phi_L$ is nonzero on an open interval, the
exponential polynomial $\sum_i c_ie^{iuz_i}$ vanishes identically. Its first
$m$ derivatives at one point form a Vandermonde system, so all $c_i=0$.
Thus the functions $\tau\mapsto\psi_L(z_i-\tau)$ are linearly independent on
every interval. The determinant above is real analytic in
$(\tau_1,\dots,\tau_m)$ and is not identically zero; a nonzero real-analytic
function cannot vanish on the open box $I_1\times\cdots\times I_m$.
\end{proof}

\begin{lemma}[Near-lattice packet Riesz bounds]
\label{lem:packet-riesz}
For lattice centres
\[
\tau_j^0=\tau_*+\frac{2\pi n_j}{L_{\rm pkt}}
\]
one has, for all large $L$,
\begin{equation}
\boxed{
2\pi\sum_j|c_j|^2
\le\left\|\sum_jc_j\psi_{L,\tau_j^0}\right\|_2^2
\le6\pi\sum_j|c_j|^2.
}
\label{eq:packet-riesz}
\end{equation}
The same inequalities, with slightly weaker fixed constants, hold when every
centre is moved by at most $T^{-100}$.
\end{lemma}

\begin{proof}
Plancherel gives
\[
\left\|\sum_jc_j\psi_{L,\tau_j^0}\right\|_2^2
=2\pi\int |\phi_L(u)|^2
\left|\sum_jc_je^{-iu\tau_j^0}\right|^2du.
\]
On the flat interval $[-L_{\rm pkt}/2,L_{\rm pkt}/2]$, the exponentials
$e^{-2\pi i n_ju/L_{\rm pkt}}$ are orthogonal; since
$|\phi_L|^2=L_{\rm pkt}^{-1}$ there, this contributes exactly
$2\pi\sum_j|c_j|^2$ and proves the lower bound. The full support has length
$L_{\rm pkt}+2\Delta<2L_{\rm pkt}$ for large $T$. By splitting it into at most
three intervals $I$ of length at most $L_{\rm pkt}$, choose for each $I$ an
interval $J\supset I$ of length exactly $L_{\rm pkt}$.  The integrand is
nonnegative, so
\[
\int_I\left|\sum_jc_je^{-2\pi i n_ju/L_{\rm pkt}}\right|^2du
\le
\int_J\left|\sum_jc_je^{-2\pi i n_ju/L_{\rm pkt}}\right|^2du
=L_{\rm pkt}\sum_j|c_j|^2.
\]
The last equality is period orthogonality on an arbitrary full-period interval;
no orthogonality on the shorter interval $I$ is asserted.  Together with
$|\phi_L|^2\le L_{\rm pkt}^{-1}$ and the factor $2\pi$ in Plancherel, summing
the at most three pieces gives the upper bound
$6\pi\sum_j|c_j|^2$.

For perturbed centres, the corresponding Gram entries depend continuously on
the centres. On the support $|u|\ll L_{\rm pkt}\asymp L$ a displacement
$T^{-100}$ changes one phase by $O(LT^{-100})$. Since the packet dimension is
$O(TL)$, the crude matrix norm perturbation is $O(T^{-99}L^2)=o(1)$, and the
fixed Riesz bounds persist. \end{proof}

\begin{lemma}[Packet dimension has logarithmic reserve]
\label{lem:packet-dimension}
For the choice \eqref{eq:packet-bandwidth-reserve}, the admissible near-lattice
packet family has
\begin{equation}
d_T=\frac{T}{2\pi}L+\frac{C_{\rm pkt}T}{2\pi}\log L+O(L).
\label{eq:packet-dimension-reserve}
\end{equation}
and therefore contains, for all large $T$, at least as many columns as there are
zeros of $F'$ in the retained dyadic rectangle even after the $O(T)$ terminal
packet collars are removed. The reserve itself is only
$O(T\log L)=o(N_T)$.
\end{lemma}

\begin{proof}
The Riemann--von Mangoldt formula gives
\[
N(T,2T)=\frac{T}{2\pi}L
+\frac{(2\log2-1)T}{2\pi}+O(\log T).
\]
By Proposition~\ref{prop:z1-zeta-decrement} and $F'=iqZ_1$ up to a nonzero
constant factor,
\[
N(F';T_1,T_2)
=\frac{T}{2\pi}L
+\frac{(2\log2-1)T}{2\pi}+O((\log T)^C).
\]
A lattice of spacing $2\pi/L_{\rm pkt}$ on a physical interval of length
$T+O(1)$ has
\[
d_T=\frac{T}{2\pi}\bigl(L+C_{\rm pkt}\log L\bigr)+O(L)
\]
columns, proving \eqref{eq:packet-dimension-reserve}. The terminal packet
collars remove only $C_{\rm del}T+o(T)$ columns for a fixed $C_{\rm del}$.
Since $C_{\rm pkt}\log L\to\infty$, the reserve term eventually dominates both
the linear Riemann--von Mangoldt secondary term and the collar loss for every
fixed $C_{\rm pkt}>0$. No further lower bound on $C_{\rm pkt}$ is required.
All later frame, mean, shell, and edge restrictions are applied
after the zero-side evaluation map has been constructed and do not enter this
surjectivity count.
\end{proof}

By Lemmas~\ref{lem:evaluation-rank}, \ref{lem:packet-riesz}, and
\ref{lem:packet-dimension}, the packet centres may therefore be chosen inside
$T^{-100}$ neighbourhoods of the lattice so that the evaluation map is
surjective while the packet synthesis retains a uniform Riesz bound.

\begin{lemma}[Coordinate restriction preserves the common proof estimates]
\label{lem:coordinate-restriction}
Let $J\subset\mathcal I_T$ be any subset of the admissible packet coordinates
and let $P_J$ be the corresponding coordinate projection. Replacing the packet
space by $P_J\mathbb C^{\mathcal I_T}$ replaces every contour/source matrix by
the compression $P_JAP_J^*$ and every positive reference Gram by
$P_JGP_J^*$. The packet Riesz lower bound is unchanged, the exact contour
identities remain valid, all edge ranks do not increase, and absolute
Hilbert--Schmidt norms of compressed remainder matrices do not increase.
Moreover the shell/frame trace-defect arguments remain valid with the retained
coordinate count in place of the full count.
\end{lemma}

\begin{proof}
Every contour, stationary, shell, and source construction is linear in the
packet synthesis coefficients, so coordinate deletion is literal matrix
compression. A subsequence of a Riesz sequence inherits the same lower and
upper Riesz constants. Rank is monotone under compression, and
$\|P_JRP_J^*\|_{S_2}\le\|R\|_{S_2}$. More importantly, if the common reference
analysis map is $W$, then on the restricted packet space it is exactly
$W_J:=WP_J^*$ and
\[
G_J=W_J^*W_J=P_JW^*WP_J^*.
\]
Since the restriction is on the input space, the least singular value of $W_J$
is no smaller than the least singular value of $W$ on the original retained
space. The primitive mixed-derivative bounds used to prove the relative $S_2$
estimates are uniform on the output fibres and are unchanged by deleting packet
inputs. Hence the relative remainder criterion
may be rerun with $W_J$ and gives the same $o(N_T)$ asymptotic (with the smaller
packet dimension). The trace-defect proofs are sums of nonnegative diagonal
contributions over packet coordinate vectors, so deleting coordinates only
removes summands and gives the same normalized estimate for the smaller
dimension. Exact algebraic contour/source identities commute with this
restriction because they hold entrywise.
\end{proof}

\begin{lemma}[Horizontal connector decay]
\label{lem:horizontal-decay}
Delete packet centres within physical distance $T/L$ of the two horizontal
heights $T_1,T_2$. This costs $O(T)=o(N_T)$ packet directions. On the
remaining packet space, every fixed curvature derivative of a horizontal
connector is $O_A(T^{-A})$ after arbitrary prescribed fixed $A$, up to fixed
polylogarithmic local-source factors.
\end{lemma}

\begin{proof}
The critical packet spacing is $\asymp L^{-1}$, so a total physical collar of
length $O(T/L)$ contains $O(T)$ centres. On either horizontal connector the
straightened packet variable has real separation at least $T/L$ from every
retained centre, while its imaginary part remains in a fixed strip. Hence
\eqref{eq:packet-decay}, after any fixed number of curvature derivatives, gives
\[
|\psi_{L,\tau}|\ll_M T^{C_M}(T/L)^{-M}.
\]
All zeta, gamma, and fixed local-source factors on the horizontal strip have at
most polynomial or polylogarithmic size. Choosing $M$ after the fixed
loss exponent gives $O_A(T^{-A})$. The same estimate holds for the adjoint
orientation and for the two mixed derivatives used in the relative
Hilbert--Schmidt criterion.
\end{proof}

For a generic $H$ with simple stationary nodes and no common zero of $H,H'$,
the residue theorem gives an exact pullback
\[
A_H=\mathcal E_H^*W_H\mathcal E_H,
\]
where $W_H$ is the direct sum of the real scalar residues and the $2\times2$
conjugate-pair blocks of Lemma~\ref{lem:stationary-blocks}. By
Lemma~\ref{lem:evaluation-rank}, $\mathcal E_H$ is surjective. Restricting to a
complement of $\ker\mathcal E_H$ and applying Sylvester's law therefore gives
\begin{equation}
\boxed{n_-(A_H)=B(H)+P(H').}
\label{eq:zero-inertia-generic}
\end{equation}

\begin{lemma}[Real zero/stationary bookkeeping]
\label{lem:real-bookkeeping}
For a real analytic function with simple real zeros and stationary points,
\[
R(H')=R(H)+2B(H)+O(1),
\label{eq:rolle-count}
\]
where $R$ denotes the number of real zeros in the interval under
consideration.
\end{lemma}

\begin{proof}
Between consecutive real zeros, Rolle's theorem forces at least one stationary
point. On every zero-free component the stationary extrema alternate. Any
additional stationary points therefore occur in pairs consisting of one good
and one bad extremum. Only the two terminal components contribute the
$O(1)$ error.
\end{proof}

For the original gauge $\mathscr Z$ and its stationary source $Z_1$, put
\[
P(s):=-\frac{\zeta'}\zeta(s),
\qquad
R(s):=\frac{Z_1(s)}{\zeta(s)}=f(s)-P(s).
\]
The functional equation gives $R(s)=-R(1-s)$. At an $m$-fold zeta zero,
\[
\zeta(s)=(s-\rho)^mu(s),\qquad u(\rho)\ne0,
\]
so
\[
Z_1(s)=m(s-\rho)^{m-1}u(\rho)+O((s-\rho)^m),
\]
and $R=Z_1/\zeta$ has a simple pole at every zeta zero, independently of its
multiplicity. Moreover, for $|t|$ large there are no zeros of $Z_1$ in
$\Re s\ge2$: there $Z_1=\zeta(f-P)$, with $P=O(1)$ and
$|f|\asymp\log|t|$. The functional equation \eqref{eq:Z1-FE} then excludes
zeros in $\Re s\le-1$. Thus every high zero of $Z_1$ relevant to the fixed
disks below lies in $-1<\Re s<2$; the corresponding assertion for the
nontrivial zeta zeros is standard from its Euler product and functional
equation.

\begin{lemma}[Local Jensen count]
\label{lem:local-z1-count}
Uniformly for $U\asymp T$,
\begin{equation}
\boxed{
N_{Z_1}([U,U+1])+N_\zeta([U,U+1])\ll\log U.
}
\label{eq:local-z1-count}
\end{equation}
\end{lemma}

\begin{proof}
Center a fixed disk at $s_U=2+iU$. Absolute convergence and Stirling give
\[
\zeta(s_U)\asymp1,
\qquad
Z_1(s_U)=\zeta(s_U)\{f(s_U)-P(s_U)\},
\qquad
|Z_1(s_U)|\gg\log U.
\]
On a fixed larger disk, the classical convexity bound for $\zeta$ and the
Cauchy estimate for $\zeta'$ give polynomial growth $U^{O(1)}$ for both
$(s-1)\zeta(s)$ and $Z_1(s)$. Jensen's formula therefore bounds the number of
zeros in the concentric smaller disk by $O(\log U)$. A fixed number of such
disks covers the strip segment of height one, proving
\eqref{eq:local-z1-count}.
\end{proof}

\begin{lemma}[Local logarithmic-derivative factorization]
\label{lem:local-log-derivative}
Let $F_0\in\{\zeta,Z_1\}$ and $U\asymp T$. Uniformly for
$-1\le\sigma\le2$ and $|t-U|\le1/10$, one has
\begin{equation}
\frac{F_0'}{F_0}(s)
=\sum_{\substack{F_0(\rho)=0\\|\rho-(2+iU)|<7/2}}
\frac1{s-\rho}+O((\log T)^{C_0}),
\label{eq:local-log-derivative}
\end{equation}
whenever $s$ is not a zero of $F_0$. The sum contains $O(\log T)$ terms,
with multiplicity.
\end{lemma}

\begin{proof}
Put $z_U=2+iU$. On the fixed disk $|s-z_U|\le5$, both $\zeta$ and $Z_1$
are holomorphic for large $U$ and have size $T^{O(1)}$. At the centre,
absolute convergence and Stirling give
\[
|\zeta(z_U)|\asymp1,
\qquad |Z_1(z_U)|\asymp\log T.
\]
Jensen's formula therefore gives $O(\log T)$ zeros, with multiplicity, in
$|s-z_U|\le4$ for either function.

Choose a radius $r_U\in[7/2,4]$ whose distance from every modulus
$|\rho-z_U|$ of these zeros is at least $c/\log T$. Such a radius exists
because there are only $O(\log T)$ forbidden intervals and the ambient radial
interval has fixed length. Factor all zeros in $|s-z_U|<r_U$:
\[
F_0(s)=e^{g_U(s)}
\prod_{|\rho-z_U|<r_U}(s-\rho)
\]
on the simply connected disk $|s-z_U|<r_U$, after removing the zeros from the
quotient. On $|s-z_U|=r_U$ the polynomial-growth upper bound for $F_0$, the
separation of the boundary from the zeros, and the $O(\log T)$ zero count give
\[
\Re g_U(s)=\log|e^{g_U(s)}|\ll \log T\log\log T.
\]
At $z_U$, the lower bounds above and
$|z_U-\rho|\le4$ give
$\Re g_U(z_U)\gg-\log T$. Applying Borel--Carath\'eodory to
$g_U-g_U(z_U)$ and then Cauchy's estimate on the fixed inner disk
$|s-z_U|\le 31/10$ yields
\[
|g_U'(s)|\ll(\log T)^{C_0}.
\]
The entire horizontal segment $-1\le\sigma\le2$, $|t-U|\le1/10$, lies in
that inner disk. Logarithmic differentiation gives the zero sum over
$|\rho-z_U|<r_U$. The additional zeros with
$7/2\le|\rho-z_U|<r_U$ are a fixed positive distance from the inner disk and
there are only $O(\log T)$ of them, so their contribution is absorbed into the
error in \eqref{eq:local-log-derivative}.
\end{proof}

\begin{proposition}[Good-height stationary-count decrement]
\label{prop:z1-zeta-decrement}
There are heights
\[
T_1\in[T,T+1],\qquad T_2\in[2T,2T+1]
\]
such that the boundary rectangle $-1\le\sigma\le2$, $T_1\le t\le T_2$
avoids all zeros and poles of $R$, and already on this good rectangle
\begin{equation}
\boxed{
N(Z_1;T_1,T_2)-N(\zeta;T_1,T_2)
=O((\log T)^C)=o(N_T).
}
\label{eq:z1-zeta-good-decrement}
\end{equation}
Returning the two horizontal heights to $T$ and $2T$ gives
\begin{equation}
\boxed{
N(Z_1;T,2T)-N(\zeta;T,2T)
=O((\log T)^C)=o(N_T).
}
\label{eq:z1-zeta-decrement}
\end{equation}
Equivalently, because $q$ is nonvanishing, the same two estimates hold with
$(Z_1,\zeta)$ replaced by $(F',F)$.
\end{proposition}

\begin{proof}
Fix $A>4$ and set $\delta_T=(\log T)^{-A}$. By
Lemma~\ref{lem:local-z1-count}, each unit interval contains only $O(\log T)$
ordinates of zeros of $\zeta Z_1$. The union of intervals of radius
$\delta_T$ around them has total length
\[
O(\delta_T\log T)=O((\log T)^{1-A})<\frac12.
\]
Hence $T_1,T_2$ can be chosen as stated, with the two horizontal sides at
distance at least $\delta_T$ from every zero or pole of $R$ in the fixed strip.

On $\Re s=2$, absolute convergence and Stirling give
\[
R(2+it)=\frac12\log\frac{t}{2\pi}+O(1),
\qquad T\le t\le2T+1.
\]
For large $T$ this lies in a fixed acute sector around the positive real axis,
so the change of argument on the right vertical side is $O(1)$. The symmetry
$R(s)=-R(1-s)$, together with
$R(\overline s)=\overline{R(s)}$, gives explicitly
\[
R(-1+it)=-R(2-it)=-\overline{R(2+it)}.
\]
Hence the left vertical image is the reflected right vertical image (with the
boundary orientation reversed), and its change of argument is also $O(1)$.

Consider one good horizontal side, at height $U\in\{T_1,T_2\}$. Since
\[
\frac{R'}R=\frac{Z_1'}{Z_1}-\frac{\zeta'}\zeta,
\]
Lemma~\ref{lem:local-log-derivative} applies separately to $Z_1$ and $\zeta$.
The selected horizontal line stays at distance at least $\delta_T$ from every
zero of either function in the fixed strip, and each local zero sum has only
$O(\log T)$ terms. Therefore, uniformly for $-1\le\sigma\le2$,
\begin{equation}
\left|\frac{R'}R(\sigma+iU)\right|
\ll (\log T)^{C_0}+\delta_T^{-1}\log T
\ll(\log T)^{A+C_1}.
\label{eq:good-horizontal-logderivative}
\end{equation}
Integrating over the bounded horizontal length gives a polylogarithmic argument
variation. Thus the total boundary variation of $\arg R$ is
$O((\log T)^C)$.

The argument principle identifies this variation divided by $2\pi$ with the
number of zeros of $R$ minus its poles. Since the divisor of $R=Z_1/\zeta$ is
exactly the divisor difference of $Z_1$ and $\zeta$, this proves
\eqref{eq:z1-zeta-good-decrement}. Returning to $T,2T$ changes
either zero count by only $O(\log T)$ by Lemma~\ref{lem:local-z1-count}, proving
\eqref{eq:z1-zeta-decrement}.
\end{proof}

Section~\ref{sec:multiplicity-endgame} transfers the generic zero-side inertia
identity \eqref{eq:zero-inertia-generic} back to the original Hardy-gauge form.
All right-edge and remainder estimates are proved for the unperturbed form
$A_F$; the generic perturbation is used only in the final finite-dimensional
inertia comparison.
